\documentclass{situationist-preprint}

\title{Husserl's Diagram of Time-Consciousness Reconsidered}
\msid{PHE-05-b}
\msfamily{PHE-05}
\msdiscipline{Phenomenology}
\mstype{Companion manuscript}
\mskeywords{time-consciousness; retention; protention; living present; diagram}

\begin{document}
\maketitle

\begin{abstract}
Husserl's lectures on internal time-consciousness contain a diagram: a horizontal line of now-points, oblique lines along which each point sinks into the past, and vertical lines representing the continuum of retentions held at one phase. We argue that the diagram, read as it is drawn, misrepresents the living present in two ways: its obliques run on without limit, and it has no protentional segment. Corrected by the retentional horizon and by protention, the diagram represents exactly the structure that PHE-05-a defends: the living present is a bounded vertical section, not a point on the horizontal line.
\end{abstract}

\section{The diagram}

In \S10 of the lectures on the phenomenology of internal time-consciousness, Husserl draws a horizontal line $OE$, the series of now-points, and from each point an oblique line descending below it, representing the sinking of that now into the past \citep{Husserl1966}. A vertical line dropped from a point $E$ cuts the obliques of the earlier points; its segment represents the continuum of retentions held at $E$, each earlier now given as just-past at a depth proportional to its distance. The vertical, not the point $E$, represents what is experienced at $E$, and this is the diagram's main lesson: the experience of succession at a phase is a continuum, not an instant.

That lesson is Theorem~4.3 of PHE-05-a: at every phase after the first, the living present is not an instant but contains other phases, ordered by depth, with its retentional extent fixed by the horizon. The diagram nonetheless suggests two things that PHE-05-a denies.

\section{Two corrections}

\emph{The obliques do not run on.} As drawn, every oblique continues downward indefinitely, and the vertical from $E$ meets every earlier oblique. Read literally, the diagram gives each phase retentions of every earlier phase, however remote. Husserl himself says that the retentional continuum fades and ends in an empty retention that no longer differentiates, and Principle~P3 of PHE-05-a states the bound: no phase is retained at a depth greater than the retentional horizon $K$. The corrected diagram cuts every oblique at depth $K$, so that the lower boundary of the region of retention is a line parallel to $OE$ at distance $K$. The vertical from $E$ then has finite length, and the living present is the bounded segment that Definition~3.5 of PHE-05-a describes. What lies below the cut is accessible only to recollection, which re-presents and does not retain.

\emph{The diagram needs an upper part.} The 1905 diagram represents only retention. Protention, which Husserl came to treat as the counterpart of retention in later manuscripts, has no place in it. A corrected diagram extends the vertical above $OE$ by a protentional segment, open and less determinate than the retentional one, as Remark~4.5 of PHE-05-a describes. The upper segment is not symmetrical with the lower one, since what is protended is anticipated in its type and not held in its content.

\section{What the diagram does not show}

One feature of the diagram is easily over-read. The horizontal line looks like a privileged present moving along a line of objective time, and the figure thus seems to depict a moving now. It does not need to be read so. The horizontal line is the order of phases of experience, and each vertical is the structure of one phase; a tenseless theory, on which no moment is metaphysically present, can accept both, as PHE-05-a argues. The diagram represents a structure of consciousness that occurs at every phase, not a property that one phase has and others lack.

\begin{objection}
Dainton argues that the retentional model, which the diagram encodes, cannot explain how a succession is experienced as a single unified episode, and that extensionalist accounts do better \citep{Dainton2000}.
\end{objection}

\begin{reply}
The corrected diagram answers the worry without becoming extensionalist. The unity of an experienced succession is the unity of one vertical segment, whose phases are given together as ordered by depth; the longitudinal intentionality of retention, which PHE-05-a uses in Proposition~5.1, gives the flow its unity across verticals.
\end{reply}

\section{Conclusion}

Husserl's diagram is right about the vertical and wrong about its length and its upper end. Bounded below by the retentional horizon and extended upward by protention, it represents the living present as PHE-05-a describes it: a structured thickness whose extent is fixed by the retentional horizon.

\begin{thebibliography}{99}\small
\bibitem[Dainton(2000)]{Dainton2000} Dainton, B. (2000). \emph{Stream of Consciousness: Unity and Continuity in Conscious Experience}. Routledge.
\bibitem[Husserl(1966)]{Husserl1966} Husserl, E. (1966). \emph{Zur Phänomenologie des inneren Zeitbewusstseins (1893--1917)}. Ed. R.~Boehm. Husserliana X. Martinus Nijhoff.
\end{thebibliography}

\end{document}
